\documentclass[conference]{IEEEtran}
% Official IEEEtran conference template (IEEEtran.cls v1.8b, from CTAN). Build: make (in doc/)
\usepackage{cite}
\usepackage{amsmath,amssymb}
\usepackage{graphicx}
\usepackage{booktabs}
\usepackage{xcolor}
\usepackage[hidelinks]{hyperref}

\graphicspath{{figures/}}

% ---------------------------------------------------------------------
% Drafting helpers. Set \draftfalse before submitting: guidance notes
% disappear, and any placeholder left behind prints in red so it is easy to spot.
\newif\ifdraft \drafttrue
\newcommand{\ph}[1]{\textcolor{red}{[#1]}}                      % placeholder to replace
\newcommand{\tbd}{\ph{--}}                                       % result not available yet
\newcommand{\guide}[1]{\ifdraft\par\noindent{\color{gray}\small\itshape #1}\par\fi} % writing guidance
% Figure placeholder: a framed box until the file exists in figures/
\newcommand{\figbox}[2][\columnwidth]{%
  \IfFileExists{figures/#2}{\includegraphics[width=#1]{#2}}%
  {\fbox{\parbox[c][3cm][c]{0.9#1}{\centering\color{red}\small Figure: \detokenize{#2}}}}}
% ---------------------------------------------------------------------

\begin{document}
\bstctlcite{IEEEexample:BSTcontrol}

\title{Does D3 Need a Pre-Trained Encoder?\\
Random-Weight Ablations of Training-Free AI-Generated Video Detection}

\author{\IEEEauthorblockN{Roland Chen}
\IEEEauthorblockA{School of Electrical Engineering and Computer Science\\
University of Ottawa, Ottawa, Canada\\
rchen095@uottawa.ca}
\and
\IEEEauthorblockN{Najwa Sougrati}
\IEEEauthorblockA{School of Electrical Engineering and Computer Science\\
University of Ottawa, Ottawa, Canada\\
nsougrat@uottawa.ca}}

\maketitle

\begin{abstract}
D3 is a training-free detector of AI-generated video. It encodes each frame
with a pre-trained visual encoder and scores a video by the volatility of the
second-order differences of its frame features. Its authors report similar
accuracy with ten encoders, but all ten are pre-trained. We ask whether D3
needs a pre-trained encoder at all. On GenVideo-Val, we compare two pre-trained
encoders (X-CLIP B/16 and ResNet-18), the same two architectures with random
weights (three seeds each), and a random projection of downsampled frames with
no network. Following a decision rule fixed before any run, the randomly
initialized encoders reach \ph{XX.X} and \ph{XX.X} mAP, against \ph{XX.X} for
pre-trained X-CLIP, so the hypothesis is \ph{accepted / rejected / partially
supported}. \ph{One sentence on the random projection.} We also control for a
frame-count confound in the evaluation protocol and report chance-corrected
margins, bootstrap confidence intervals and paired tests. \ph{One sentence on
what this means for D3.}
\end{abstract}

\begin{IEEEkeywords}
AI-generated video detection, training-free detection, random features,
second-order features, ablation study
\end{IEEEkeywords}

% =====================================================================
\section{Introduction}
\label{sec:intro}

Generated video is becoming hard to tell apart from real footage, and
detectors trained on one generator often fail on the next. D3~\cite{zheng2025d3}
avoids training altogether. It passes each frame through a pre-trained visual
encoder, measures how far the features move between consecutive frames (first
order), then how that movement changes (second order). The score of a video is
the standard deviation of the second-order sequence. With an X-CLIP
encoder~\cite{ni2022xclip}, D3 reaches 97.72 mAP on
GenVideo-Val~\cite{chen2024demamba}.

The D3 authors compare ten encoders and report between 94.28 and 97.72 mAP. They conclude that second-order features stay meaningful
across feature spaces~\cite[Sec.~5.1]{zheng2025d3}. However, all ten encoders
are pre-trained. If \emph{which} pre-trained encoder is used barely matters,
the next question is whether pre-training is needed at all. Networks with
random weights are known to give useful features~\cite{jarrett2009,saxe2011},
and random projections preserve distances~\cite{johnson1984}. If D3 relies
mainly on its temporal statistic, an untrained encoder should do almost as
well.

We test this with five inference-only conditions: two reproductions with
pre-trained encoders, the same two architectures with random weights, and a
random projection with no network. The hypothesis and the decision thresholds
were fixed in our proposal before any experiment. Our contributions are:
\begin{itemize}
  \item a reproduction of two encoder conditions from the D3 paper;
  \item an ablation with randomly initialized encoders (three seeds each) and
        a no-network random projection;
  \item a control for a frame-count confound that we found in the evaluation
        protocol;
  \item a statistical reading of the pre-registered thresholds: margins over
        chance, bootstrap confidence intervals and paired tests.
\end{itemize}
Our code is publicly available (see the Code Availability section).

% =====================================================================
\section{Related Work}
\label{sec:related}

\textbf{AI-generated video detection.} Most detectors are trained on generated
videos. Chen \emph{et al.}~\cite{chen2024demamba} release GenVideo, a
million-scale benchmark, and its validation split GenVideo-Val, which we use.
D3~\cite{zheng2025d3} needs no training: it only uses a frozen encoder and a
fixed statistic.

\textbf{Random and untrained networks as feature extractors.} Jarrett
\emph{et al.}~\cite{jarrett2009} found that a two-stage convolutional network
with random filters reaches almost 63\% accuracy on Caltech-101, provided
that it has the right non-linearities and pooling. Saxe
\emph{et al.}~\cite{saxe2011} showed that convolutional pooling architectures
are frequency selective and translation invariant even with random weights,
and used random weights to rank architectures without training. Random Fourier
features~\cite{rahimi2007} and extreme learning machines~\cite{huang2006elm}
fix a random first layer and learn only a linear output. Giryes
\emph{et al.}~\cite{giryes2016} showed that networks with random Gaussian
weights preserve the metric structure of their input. Deep Image
Prior~\cite{ulyanov2018dip} showed that the structure of an untrained
convolutional network already captures much of the statistics of natural
images.

\textbf{Pre-training without natural images.} Kataoka
\emph{et al.}~\cite{kataoka2020fractal} pre-train networks on fractal images,
and Baradad \emph{et al.}~\cite{baradad2021noise} on synthetic noise
processes. Both obtain useful features without natural images. We go one step
further and use no training at all.

\textbf{Random projections.} The Johnson--Lindenstrauss
lemma~\cite{johnson1984} states that a random linear map to a lower dimension
preserves pairwise distances up to a small relative error, and
Achlioptas~\cite{achlioptas2003} showed that simple $\pm1$ matrices are
enough. This motivates our no-network baseline.

% =====================================================================
\section{Method}
\label{sec:method}

\subsection{The D3 Score}

Let $x_1,\dots,x_T$ be the frames of a video and $E$ an encoder, with
features $f_t = E(x_t) \in \mathbb{R}^d$. Following the official
implementation~\cite{d3code}, D3 computes
\begin{align}
  d_t &= \lVert f_{t+1} - f_t \rVert_2, & t &= 1,\dots,T-1,\\
  a_t &= d_{t+1} - d_t,                 & t &= 1,\dots,T-2,\\
  s   &= \operatorname{std}(a_1,\dots,a_{T-2}), &&
\end{align}
where $\operatorname{std}$ is the sample standard deviation. In the official
evaluation, a \emph{higher} score indicates a \emph{real} video: average
precision is computed with real videos as the positive class. D3 has no
trainable parameter, so changing $E$ is the only way to change the detector.

\subsection{Experimental Conditions}

Table~\ref{tab:conditions} lists the five conditions. Condition~1 is the
paper's main setting: the vision encoder of X-CLIP B/16~\cite{ni2022xclip}, a
ViT-B/16~\cite{dosovitskiy2021vit} with cross-frame communication, using the
pooled output. Condition~2 is ResNet-18~\cite{he2016resnet} pre-trained on
ImageNet, without its classification layer. Conditions~3 and~4 use the same
architectures with random weights: torchvision's default initialization for
ResNet-18, and the Hugging Face default initialization for X-CLIP (zero-mean
Gaussian weights). Condition~5 has no network. It average-pools each frame to
$32\times32$ pixels and multiplies the resulting 3072-dimensional vector by a
fixed matrix $W \in \mathbb{R}^{3072\times768}$ with entries drawn from
$\mathcal{N}(0, 1/768)$. The seed fixes the random weights. The same seed
always gives the same weights.

\begin{table}[t]
\caption{Experimental conditions. Conditions 3--5 are run with seeds 0, 1
and 2.}
\label{tab:conditions}
\centering
\footnotesize
\setlength{\tabcolsep}{4pt}
\begin{tabular}{@{}clll@{}}
\toprule
\# & Encoder & Weights & Dim. \\
\midrule
1 & X-CLIP B/16 vision encoder & pre-trained (Kinetics-400) & 768 \\
2 & ResNet-18                  & pre-trained (ImageNet)     & 512 \\
3 & ResNet-18                  & random                     & 512 \\
4 & X-CLIP B/16 vision encoder & random                     & 768 \\
5 & Random projection          & random, no network         & 768 \\
\bottomrule
\end{tabular}
\end{table}

\subsection{Hypothesis and Decision Rule}
\label{sec:rule}

\textbf{Hypothesis.} D3's performance does not depend on the learned content
of its encoder: with random weights, D3 keeps most of its performance.

\textbf{Decision rule (fixed before any run).} We use the mean mAP over three
seeds of conditions~3 and~4. If both exceed 88, the hypothesis is accepted. If
both are below 70, it is rejected. Any other outcome, including disagreement
between the two architectures, indicates partial dependence. Condition~5 is
read separately: reaching 88 with no network would show that no learned
representation is needed at all. Conditions~1 and~2 are reproductions,
checked against the reported 97.72 and 96.39 mAP.

\textbf{Reading the thresholds.} The thresholds were set before any result,
and we do not move them. Instead, we add statistics beside them. First, chance
performance is not 50 mAP. With real videos as the positive class, a random
score has an expected AP equal to the share of real videos in the run. This
share varies with the size of each fake subset, from 50.0\% to 94.7\%
(Table~\ref{tab:subsets}). The mean over the ten subsets, $m_0 = 60.23$,
depends only on the class counts. We measure each condition $c$ by the share
of the margin over chance that it keeps:
\begin{equation}
  \rho_c = \frac{\mathrm{mAP}_c - m_0}{m_{\mathrm{ref}} - m_0},
\end{equation}
with $m_{\mathrm{ref}} = 97.72$, the X-CLIP result reported in the paper.
Since $m_0$ and $m_{\mathrm{ref}}$ are both known before any run, the
thresholds translate into fixed shares: 88 mAP means keeping 74.1\% of the
margin over chance, and 70 mAP means keeping 26.1\%. Accepting the hypothesis
thus requires random encoders to keep about three quarters of what the
pre-trained encoder gains over chance, and rejecting it requires them to lose
about three quarters. Second, we attach 95\% bootstrap confidence intervals to
every mAP (Sec.~\ref{sec:stats}). We call a verdict \emph{robust} when the
interval of each random encoder lies entirely on the same side of the
threshold as its mean.

% =====================================================================
\section{Experimental Setup}
\label{sec:setup}

\subsection{Data}

GenVideo-Val~\cite{chen2024demamba} contains real videos from
MSR-VTT~\cite{xu2016msrvtt} and ten subsets of generated videos
(Table~\ref{tab:subsets}). Each run pairs the same 1000 real videos with one
fake subset, capped at its first 1000 videos. Videos with fewer than 8 frames
cannot be scored and are excluded: 31 fakes in total (2 in Gen2, 1 in
MoonValley and 28 in WildScrape). mAP is the unweighted mean of the ten
per-subset APs, as in the paper, so Sora (56 fakes) counts as much as Gen2
(1000 fakes).

\begin{table}[t]
\caption{Fake subsets of GenVideo-Val and AP per condition (paper protocol,
seeds averaged). ``8-fr.'': share of fakes scored on 8 frames instead of 16;
all real videos are scored on 16. ``Chance'': expected AP of a random score.}
\label{tab:subsets}
\centering
\footnotesize
\setlength{\tabcolsep}{2.5pt}
\begin{tabular}{@{}lrrr|rrrrr@{}}
\toprule
 & & 8-fr. & & \multicolumn{5}{c}{AP (\%) for condition} \\
Subset & Fakes & (\%) & Chance & 1 & 2 & 3 & 4 & 5 \\
\midrule
Crafter     & 1000 &  86.0 & 50.00 & \tbd & \tbd & \tbd & \tbd & \tbd \\
Gen2        & 1000 &   0.1 & 50.00 & \tbd & \tbd & \tbd & \tbd & \tbd \\
HotShot     &  607 & 100.0 & 62.23 & \tbd & \tbd & \tbd & \tbd & \tbd \\
Lavie       & 1000 &   0.1 & 50.00 & \tbd & \tbd & \tbd & \tbd & \tbd \\
ModelScope  &  700 &   0.0 & 58.82 & \tbd & \tbd & \tbd & \tbd & \tbd \\
MoonValley  &  581 & 100.0 & 63.25 & \tbd & \tbd & \tbd & \tbd & \tbd \\
MorphStudio &  619 &   0.0 & 61.77 & \tbd & \tbd & \tbd & \tbd & \tbd \\
Show\_1     &  700 &   0.0 & 58.82 & \tbd & \tbd & \tbd & \tbd & \tbd \\
Sora        &   56 &   0.0 & 94.70 & \tbd & \tbd & \tbd & \tbd & \tbd \\
WildScrape  &  896 &  12.3 & 52.74 & \tbd & \tbd & \tbd & \tbd & \tbd \\
\midrule
mAP         &      &       & 60.23 & \tbd & \tbd & \tbd & \tbd & \tbd \\
\bottomrule
\end{tabular}
\end{table}

\subsection{Preprocessing}

We keep the official preprocessing~\cite{d3code}, so that conditions~1 and~2
reproduce the paper. Frames are extracted at 8~fps from a 3-second clip. For
videos longer than 4~s, the clip starts at a random whole second. This applies
to all real videos, all Sora videos and 212 WildScrape videos. Each video is
scored on its first 16 frames (2~s), or on its first 8 frames if it has fewer
than 16. Each frame is cropped by 10\% on both sides of its longer dimension,
resized to $224\times224$, and normalized with ImageNet statistics. The same
normalization is used for every encoder, including X-CLIP. As in the original
code, frames are read in OpenCV's BGR channel order. X-CLIP's vision encoder
exchanges information across frames in groups of 8, so a 16-frame video is
processed as two independent groups.

\subsection{Frame-Count Control}
\label{sec:f8}

Under this protocol, every real video is scored on 16 frames, but every fake
in HotShot and MoonValley, and 86\% of the fakes in Crafter, are scored on 8
(Table~\ref{tab:subsets}). With 8 frames, the score is the standard deviation
of 6 values over 1~s instead of 14 values over 2~s. Frame count alone would
therefore separate real from fake perfectly on HotShot and MoonValley. To
measure how much of each result this explains, we repeat every condition with
all videos scored on their first 8 frames (seed 0 only).

\subsection{Statistical Analysis}
\label{sec:stats}

We save the score of every video, so all statistics are computed from the same
runs. We estimate 95\% percentile bootstrap confidence
intervals~\cite{efron1993} with $B = \ph{2000}$ replicates. Each replicate
resamples the 1000 real videos once, since they are shared by all subsets,
and the fakes of each subset independently. It then recomputes every AP and
the mAP. For random conditions, the three seeds are averaged within each
replicate. To compare a pre-trained and a random encoder, we use the same
resampled videos for both and report the confidence interval of the mAP
difference. We also report AUC, whose chance level is 0.5 whatever the class
balance, and the Spearman correlation of per-video scores between encoders.
\guide{Update this paragraph if the analysis script differs.}

\subsection{Implementation}

All 160 runs (110 for the paper protocol and 50 for the frame-count control)
were made on one machine, with an AMD Ryzen AI 7 350 CPU and a Radeon 860M GPU
(ROCm), using PyTorch 2.8, torchvision 0.23 and Hugging Face Transformers 4.57.
A full run took \ph{time}.

% =====================================================================
\section{Results}
\label{sec:results}

\guide{Fill in once Najwa's runs are in. Report numbers with two decimals in
the tables and one in the text.}

\subsection{Reproduction}

\ph{Pre-trained X-CLIP reaches XX.XX mAP and pre-trained ResNet-18 XX.XX,
against 97.72 and 96.39 reported~\cite{zheng2025d3}. One sentence on any gap
and its likely cause: video selection, random clip choice, excluded short
videos.}

\subsection{Random Encoders and the Pre-Registered Verdict}

\begin{table}[t]
\caption{mAP per condition. Mean $\pm$ standard deviation over three seeds for
random conditions; 95\% bootstrap CI; $\rho$: share of the margin over chance
kept (Sec.~\ref{sec:rule}); last column: all videos scored on 8 frames.}
\label{tab:summary}
\centering
\footnotesize
\setlength{\tabcolsep}{3pt}
\begin{tabular}{@{}clcccc@{}}
\toprule
\# & Condition & mAP & 95\% CI & $\rho$ (\%) & mAP, 8 fr. \\
\midrule
1 & X-CLIP, pre-trained  & \tbd         & \tbd & \tbd & \tbd \\
2 & ResNet-18, pre-trained & \tbd       & \tbd & \tbd & \tbd \\
3 & ResNet-18, random    & \tbd$\pm$\tbd & \tbd & \tbd & \tbd \\
4 & X-CLIP, random       & \tbd$\pm$\tbd & \tbd & \tbd & \tbd \\
5 & Random projection    & \tbd$\pm$\tbd & \tbd & \tbd & \tbd \\
\midrule
  & Chance               & 60.23        & --   & 0    & 60.23 \\
\bottomrule
\end{tabular}
\end{table}

\ph{Verdict, stated exactly as the pre-registered rule gives it:
``ResNet-18-rand reaches XX.X and X-CLIP-rand XX.X mAP; both are above 88, so
the hypothesis is accepted.'' Then: is it robust (bootstrap CIs vs
thresholds)? How large is the paired difference to the pre-trained encoder of
the same architecture, with its CI? Seed spread. RandProj, read separately.}

\subsection{Frame-Count Control}

\ph{mAP with all videos on 8 frames (Table~\ref{tab:summary}, last column), and
the change on HotShot, MoonValley and Crafter, the three subsets where frame
count differs. Does the verdict change under the control?}

\begin{figure}[t]
\centering
\figbox{f8_per_subset.pdf}
\caption{\ph{AP per subset under the paper protocol and with all videos scored
on 8 frames, for each condition.}}
\label{fig:f8}
\end{figure}
\guide{Optional figure: drop it if the 5-page limit is tight.}

\subsection{Where Pre-Training Matters}

\ph{Subsets with the largest gap between pre-trained and random encoders
(Table~\ref{tab:subsets}), with AUC. Spearman correlation of per-video scores
between encoders: do random encoders rank videos like the pre-trained ones?}

% =====================================================================
\section{Discussion}
\label{sec:discussion}

\subsection{What Carries the Performance}

\ph{Interpretation, depending on the verdict: is the second-order statistic,
the architecture, or the learned representation doing the work?}

For the random projection, a possible explanation follows from the
Johnson--Lindenstrauss lemma~\cite{johnson1984,achlioptas2003}. For a fixed
vector $v$, $\lVert W^\top v\rVert^2 / \lVert v\rVert^2$ follows a $\chi^2$
distribution with $k = 768$ degrees of freedom divided by $k$. Its relative
standard deviation is $\sqrt{2/k} \approx 5\%$. The distances $d_t$ that D3
computes on projected features are therefore close to the distances between
the raw $32\times32$ frames, whatever the seed. Condition~5 then measures D3
on downsampled pixels, and its result should depend little on the seed:
\ph{observed seed spread}.

\subsection{Implications for D3}

\ph{What the result means for the claim that second-order features remain
meaningful across feature spaces~\cite[Sec.~5.1]{zheng2025d3}: does it extend
to untrained feature spaces, or is it narrower than stated? Practical
consequence: cost, or the choice of encoder.}

\subsection{Limitations}

\begin{itemize}
  \item We test one benchmark, GenVideo-Val. The other datasets of the
        paper (EvalCrafter, VideoPhy, VidProM) are not tested.
  \item Every run uses the same 1000 real videos, so errors on the real side
        are shared across subsets. Subsets are capped at 1000 fakes.
  \item mAP weights subsets equally. Sora has only 56 fakes and a chance AP of
        94.7, so it carries little information about any condition.
  \item Each real video is scored on one random 2-second clip. Our seeds do
        not measure how much the choice of clip changes the results.
  \item Random conditions use three seeds, and the frame-count control uses
        one.
  \item We keep the original preprocessing (ImageNet normalization and BGR
        order for every encoder). It may not suit X-CLIP, which was trained
        with CLIP statistics on RGB frames.
  \item Random weights follow library defaults. In inference mode, the batch
        normalization layers of the random ResNet-18 use their initial
        statistics and act as the identity. Other initializations may behave
        differently.
\end{itemize}

% =====================================================================
\section{Conclusion}
\label{sec:conclusion}

\ph{Three to four sentences: the question, the pre-registered verdict and its
robustness, the frame-count control, and what it means for training-free
detection.}

% =====================================================================
\section*{Code Availability}

Our code is available at
\url{https://github.com/Lylichen07/D3-Ablation-Review}
\ph{check that the repository is public, or state where the code is
submitted}. It builds on the official D3 implementation~\cite{d3code} (MIT
License). The D3 score and the preprocessing are unchanged. Data
preparation follows DeMamba~\cite{chen2024demamba}. Our additions are the
randomly initialized encoders, the random-projection baseline, the seeding of
random weights, the exclusion of videos with fewer than 8 frames, per-video
score files, the 8-frame control, and the scripts that run the experiments and
analyze the results.

{\small
\bibliographystyle{IEEEtran}
\bibliography{references}
}

\end{document}
